% !PRINTED 672
\weberpage{672}
%<<CONT
2. \gesperrt{If two $n$-gons are equivalent to a third, then they are also equivalent to one another. Accordingly one unites equivalent polygons $\mathfrak{A}, \mathfrak{A}', \ldots$ into a polygon class $A$. Every polygon is contained in one class, and in only one class.}

There also exist polygons that are equivalent to no other, and that therefore each form a separate class by themselves. These are called \gesperrt{isolated polygons}.

3. \gesperrt{If $\mathfrak{A} \sim \mathfrak{A}'$, $\mathfrak{B} \sim \mathfrak{B}'$, then $\mathfrak{A}\mathfrak{B} \sim \mathfrak{A}'\mathfrak{B}'$.}

This follows from

\weq{4}{\eta = \dfrac{\mathfrak{A}}{\mathfrak{A}'}\,, \quad \eta_1 = \dfrac{\mathfrak{B}}{\mathfrak{B}'}\,, \quad \eta\eta_1 = \dfrac{\mathfrak{A}\mathfrak{B}}{\mathfrak{A}'\mathfrak{B}'}.}

The class $C$ to which the product $\mathfrak{A}\mathfrak{B}$ of the classes $A, B$ belongs therefore contains all products of a polygon of the class $A$ with a polygon of the class $B$ (in certain circumstances still others as well), and is therefore completely determined by the classes $A$ and $B$. On this rests the multiplication (composition) of the classes, which is expressed in the formula:

\weq{5}{C = AB = BA.}

From an equation between three classes $A, B, M$

\[
MA = MB
\]

it follows accordingly that $A = B$. For if $\mathfrak{M}\mathfrak{A} \sim \mathfrak{M}\mathfrak{B}$, then by (2) it follows that $\mathfrak{A} \sim \mathfrak{B}$. If $\mathfrak{A} \sim \mathfrak{A}'$ and

\[
\mathfrak{A}\mathfrak{B} = \mathfrak{C} \quad \text{and} \quad \mathfrak{A}'\mathfrak{B} = \mathfrak{C}'\,,
\]

then $\mathfrak{C} \sim \mathfrak{C}'$ as well. If, then, $\mathfrak{A}$ is contained in $\mathfrak{C}$, every polygon equivalent to $\mathfrak{A}$ is contained in a polygon $\mathfrak{C}'$ equivalent to $\mathfrak{C}$.

4. \gesperrt{We call a class $C$ divisible by a class $A$ if a polygon of $A$ is contained in a polygon of $C$, and we set}

\weq{6}{AB = C.}

\gesperrt{The class $B$ is completely determined by $A$ and $C$.}

\begin{center}
\textbf{§ 192. Families of Polygons.}
\end{center}

If $\mathfrak{A}, \mathfrak{A}_1, \mathfrak{A}_2, \ldots, \mathfrak{A}_s$ are polygons of a class $A$, then there are $s$ functions in $\Omega$:

\weq{1}{\eta_1 = \dfrac{\mathfrak{A}_1}{\mathfrak{A}}\,, \quad \eta_2 = \dfrac{\mathfrak{A}_2}{\mathfrak{A}}\,, \quad \cdots, \quad \eta_s = \dfrac{\mathfrak{A}_s}{\mathfrak{A}}\,.}
